\documentclass[a4paper,UTF8,12pt]{ctexart}
\usepackage{Mystyle}
\title{数学基础}
\author{C\'elian Sun}
\date{\quad }
\begin{document}
\thispagestyle{empty}
\maketitle
\tableofcontents
\newpage
\color{ZhengWen}
% \input{Section/Set_Axiomatics.tex}
% \input{Section/Classical_Calculus.tex}
% \input{Section/Complex_Analysis.tex}
% \input{Section/Real_Analysis.tex}
% \input{Section/Function_Analysis.tex}
% \input{Section/Foutier_Analysis.tex}
% \input{Section/Topology_Geometry.tex}
% \input{Section/Algebra_NumberTheory.tex}
% \input{Section/Finite_Element.tex}
%\input{Section/Machine_Learning.tex}
%\input{Section/Numerical_Analysis.tex}

\newpage
\section{两个求和公式及应用}
\subsection{Euler-Maclaurin求和公式}
级数的一个重要问题在于如何判断其收敛性,这一问题本质上就是分析级数的渐近性质,也就是怎么增长或衰减的.这在数论中(尤其是解析数论里)十分重要\cite{apostol1998introduction}.下面给出一般的Euler-Maclaurin求和公式\cite{enwiki:1351518612},它是级数收敛性判断的一个重要工具.
\begin{Theorem}{Euler-Maclaurin求和公式}
设$a,b\in\mathbb{Z}$且$a<b$,函数$f\in C^{m}([a,b])$,$m$为正整数.则
$$
\sum_{n=a}^{b}f(n)=\int_{a}^{b}f(x)dx+\frac{f(a)+f(b)}{2}+\sum_{k=1}^{\lfloor m/2\rfloor }\frac{B_{2k}}{(2k)!}(f^{(2k-1)}(b)-f^{(2k-1)}(a))+R_m
$$
余项$R_m=\frac{(-1)^{m+1}}{m!}\int_{a}^{b}B_{m}(x-\lfloor x \rfloor)f^{(m+1)}(x)dx$.其中$B_{k}$为第$k$个Bernoulli数,$B_{m}(x)$为第$m$个Bernoulli多项式.
\end{Theorem}
这个公式并不证明.首先说一下什么是Bernoulli数和Bernoulli多项式.Bernoulli多项式是通过下面的方式定义的:
$$
B_0(x)=1, \quad B_n'(x)=nB_{n-1}(x), \quad \int_0^1 B_n(x) dx=0
$$
我们可以给出前几个Bernoulli多项式的表达式:
$$
\begin{aligned}
B_0(x) &= 1 \\
B_1(x) &= x - \frac{1}{2} \\
B_2(x) &= x^2 - x + \frac{1}{6} \\
B_3(x) &= x^3 - \frac{3}{2}x^2 + \frac{1}{2}x \\
B_4(x) &= x^4 - 2x^3 + x^2 - \frac{1}{30} \\
\end{aligned}
$$
而Bernoulli数则是Bernoulli多项式在$x=0$处的值,即$B_n = B_n(0)$.前几个Bernoulli数的值如下:
$$
B_0=1, \quad B_1=-\frac{1}{2}, \quad B_2=\frac{1}{6}, \quad B_3=0, \quad B_4=-\frac{1}{30}, \quad B_5=0, \quad B_6=\frac{1}{42}, \quad \cdots
$$
可以发现$B_{2k+1}=0$对于所有$k\geq 1$成立.这解释了为什么在Euler-Maclaurin求和公式中,除了$B_1$以外,只有偶数项出现.\par
现在我们考虑一个简单版本的Euler-Maclaurin公式:
\begin{Theorem}{}
设$a,b\in\mathbb{Z}$且$a<b$,函数$f\in C^{1}([a,b])$.则
$$
\sum_{n=a}^{b}f(n)=\int_{a}^{b}f(x)dx+\frac{f(a)+f(b)}{2}+\int_{a}^{b}(x-\lfloor x \rfloor-\frac{1}{2})f'(x)dx
$$
\end{Theorem}
\begin{Proof}
    直接计算得到
    $$
    \begin{aligned}
        \int_{a}^{b}(x-\lfloor x \rfloor-\frac{1}{2})f'(x)dx &= \sum_{k=a}^{b-1}\int_{k}^{k+1}(x-\lfloor x \rfloor-\frac{1}{2})f'(x)dx \\
            &= \sum_{k=a}^{b-1}\int_{0}^{1}(x-\frac{1}{2})f'(k+x)dx \\
            &= \sum_{k=a}^{b-1}\left[\frac{f(k)+f(k+1)}{2}-\int_{k}^{k+1}f(x)dx\right] \\
            &= \sum_{k=a}^{b}f(k)-\int_{a}^{b}f(x)dx-\frac{f(a)+f(b)}{2}
    \end{aligned}
    $$
    这就证明了结论.
\end{Proof}
我们可以考虑几个例子.首先最经典的例子是调和级数.为此我们考虑函数$f(x)=\frac{1}{x}$,则
$$
\sum_{n=1}^{N}\frac{1}{n} = \ln N+\frac{1}{2}+\frac{1}{2N}-\int_{1}^{N}\frac{x-\lfloor x \rfloor-1/2}{x^2}dx
$$
我们注意到$|x-\lfloor x \rfloor-1/2| \leqslant 1/2$,所以,当$N\to\infty$时,右侧积分$\int_{1}^{\infty}\frac{x-\lfloor x \rfloor-1/2}{x^2}dx$绝对收敛.从这里我们就可以看出来调和级数的发散性,甚至可以看出来当$N$越来越大时,调和级数的增长速度大约是$\ln N$.\par
接下来,我们考虑作业中的一个问题,即级数$\sum_{n=1}^{\infty}\frac{\ln n}{n^2+1}$.为此,计算
$$
\sum_{n=1}^{N}\frac{\ln n}{n^2+1} = \int_{1}^{N}\frac{\ln x}{x^2+1}dx+\frac{\ln N}{2(N^2+1)}+\int_{1}^{N}(x-\lfloor x \rfloor-\frac{1}{2})\left(\frac{1+x^2(1-2\ln x)}{x(x^2+1)^2}\right)dx
$$
如此一来我们要分析下面两个积分
$$
\int_{1}^{\infty}\frac{\ln x}{x^2+1}dx,\quad \int_{1}^{\infty}(x-\lfloor x \rfloor-\frac{1}{2})\left(\frac{1+x^2(1-2\ln x)}{x(x^2+1)^2}\right)dx
$$
这两个积分显然是收敛的.因此,我们就可以得出结论,级数$\sum_{n=1}^{\infty}\frac{\ln n}{n^2+1}$是收敛的.\par
你可能觉得,本来就是要判断级数的收敛性,现在又要判断积分的收敛性,这似乎没有简化问题嘛.\par
这触碰到了Euler-Maclaurin求和公式的核心,这个公式从来不是为了精确求解,而是为了近似计算,准确来讲,是为了帮助我们分析渐近行为的.\par
比如,我们熟悉的Stirling公式,我们可以用Euler-Maclaurin公式非常简单的证明出来(这里的简单是指只需要机械计算,而不用其他技巧).我们考虑$f(x)=\ln x$,则经过计算得到
$$
\begin{aligned}
    \ln N! = \sum_{n=1}^{N}\ln n &= \int_{1}^{N}\ln x dx + \frac{\ln N}{2} + \int_{1}^{N}(x-\lfloor x \rfloor-\frac{1}{2})\frac{1}{x}dx\\
        &= N\ln N - N + 1 + \frac{\ln N}{2} + \int_{1}^{N}(x-\lfloor x \rfloor-\frac{1}{2})\frac{1}{x}dx
\end{aligned}
$$
你可能发现后面的那个积分不好处理,但是根据Bernoulli多项式的性质有$B'_2(x)=2B_1(x)=2(x-\frac{1}{2})$,因此通过分部积分有
$$
\begin{aligned}
    \int_{1}^{N}(x-\lfloor x \rfloor-\frac{1}{2})\frac{1}{x}dx &= \frac{1}{2}\int_1^N \frac{1}{x}\mathrm{d}B_2(x-\lfloor x \rfloor) \\
    &= \frac{1}{2}\left\{\left[\frac{B_2(x-\lfloor x \rfloor)}{x}\right]_{1}^{N}-\int_1^N B_2(x-\lfloor x \rfloor)\mathrm{d}\left(\frac{1}{x}\right)\right\} \\
\end{aligned}
$$
对于$x$是整数而言$B_2(x-\lfloor x \rfloor)=B_2(0)=\frac{1}{6}$,因此
$$
 \ln N!= N\ln N - N + \frac{\ln N}{2} +\left(1-\frac{1}{12}\right)+\frac{1}{12N} + \frac{1}{2}\int_1^N B_2(x-\lfloor x \rfloor)\frac{1}{x^2}dx
$$
注意到$|B_2(x-\lfloor x \rfloor)| \leqslant \frac{1}{6}$,因此$\int_1^\infty B_2(x-\lfloor x \rfloor)\frac{1}{x^2}dx$是绝对收敛的.根据积分区间的可加性,有
$$
\int_1^N B_2(x-\lfloor x \rfloor)\frac{1}{x^2}dx=\int_1^\infty B_2(x-\lfloor x \rfloor)\frac{1}{x^2}dx - \int_N^\infty B_2(x-\lfloor x \rfloor)\frac{1}{x^2}dx
$$
我们记常数$C=1-\frac{1}{12}+\frac{1}{2}\int_1^\infty B_2(x-\lfloor x \rfloor)\frac{1}{x^2}dx$,此外利用$|B_2(x-\lfloor x \rfloor)| \leqslant \frac{1}{6}$得到
$$
\left|\int_N^\infty B_2(x-\lfloor x \rfloor)\frac{1}{x^2}dx\right| \leqslant \frac{1}{6} \int_N^\infty \frac{1}{x^2}dx = \frac{1}{6N}
$$
因此
$$
\ln N! = N\ln N - N + \frac{\ln N}{2} + C + O\left(\frac{1}{N}\right)
$$
两边取指数,就得到了$N! \sim C \sqrt{N}\left(\frac{N}{\mathrm{e}}\right)^N$.\par
\subsection{Abel求和公式}
Euler-Maclaurin求和公式是依赖于函数的光滑性,对于那些离散的,甚至有些带权重的级数,例如$\sum a_nf(n)$该怎么解决呢?,我们的想法是将导数转移到$f$的身上.各位已经学过Abel分部求和公式,但是我们换一个版本\cite{enwiki:1321566893}.\par
\begin{Theorem}{Abel求和公式}
设$A(x)=\sum_{1\leqslant n\leqslant x}a_n$,$f\in C^1([1,x])$,则有
$$
    \sum_{1\leqslant n\leqslant x}a_nf(n) = A(x)f(x) - \int_1^x A(t)f'(t)dt
$$
\end{Theorem}
\begin{Proof}
    固定$x$使得$k\leqslant x\leqslant k+1$.利用Abel分部求和公式有
    $$
    \begin{aligned}
        \sum_{1\leqslant n\leqslant x}a_nf(n) &= \sum_{1\leqslant n\leqslant k}a_nf(n)\\
        &= A(k)f(k)-\sum_{1\leqslant n\leqslant k-1}A(n)(f(n+1)-f(n)) \\
    \end{aligned}
    $$
    注意到$A(t)=A(n),n\leqslant t<n+1$,计算右边和式为
    $$
    \begin{aligned}
        \sum_{1\leqslant n\leqslant k-1}A(n)(f(n+1)-f(n)) &= \sum_{1\leqslant n\leqslant k-1}A(n)\int_n^{n+1}f'(t)dt \\
        &= \sum_{1\leqslant n\leqslant k-1}\int_n^{n+1}A(t)f'(t)dt\\
        &= \int_1^k A(t)f'(t)dt=\int_1^x A(t)f'(t)dt-\int_k^x A(t)f'(t)dt
    \end{aligned}
    $$
    又因为
    $$
    \int_k^x A(t)f'(t)dt= A(k)\int_k^x f'(t)dt=A(x)f(x)-A(k)f(k)
    $$
    整理上面的结果得到
    $$
    \sum_{1\leqslant n\leqslant x}a_nf(n) = A(x)f(x) - \int_1^x A(t)f'(t)dt
    $$
\end{Proof}
当$a_n=1$时,$A(x)=\lfloor x \rfloor$,这样就得到了最简单形式的Euler-Maclaurin求和公式.但我们的目的不是这个,我们将用这个公式来分析所有素数的倒数和$\sum_{p\leqslant x}\frac{1}{p}$的性质,即Mertens\texttt{'}Second Theorem\cite{enwiki:1352898521}\cite{villarino2005mertensproofmertenstheorem}.\par
我们设$a_n$为素数的示性函数,即
$$
a_n=\mathbf{1}_p(n)=
\begin{cases}
1, & \text{如果 $n$ 是素数} \\
0, & \text{其他}
\end{cases}
$$
那么$A(x)=\sum_{n\leqslant x}a_n$其实就是表示小于等于$x$的素数的个数.在数论中有专门的函数,即$\pi(x)$.取$f(x)=\frac{1}{x}$,利用Abel求和公式有
$$
\sum_{p\leqslant x}\frac{1}{p} = \frac{\pi(x)}{x} + \int_2^x \frac{\pi(t)}{t^2}dt
$$
利用Prime number theorem\cite{terencetao_zeta}有$\lim_{x \to \infty} \frac{\pi(x)}{x/\ln x} = 1$,定义误差函数$E(x)=\pi(x)-\frac{x}{\ln x}$,则存在常数$C$使得$|E(x)| \leqslant C\frac{x}{\ln^2 x}$对于所有$x\geqslant 2$成立.因此
$$
    \frac{\pi(x)}{x}= \frac{1}{\ln x} + \frac{E(x)}{x}= \frac{1}{\ln x} + O\left(\frac{1}{\ln^2 x}\right)
$$
第二个等号是因为$\left|\frac{E(x)}{x}\right|\leqslant C\frac{1}{\ln^2 x}$.接下来是关于积分的处理.首先
$$
\int_2^x \frac{\pi(t)}{t^2}dt= \int_2^x \frac{1}{t\ln t}dt + \int_2^x \frac{E(t)}{t^2}dt=\ln\ln x-\ln\ln 2 + \int_2^x \frac{E(t)}{t^2}dt
$$
利用$\left|\frac{E(t)}{t^2}\right| \leqslant C\frac{1}{t\ln^2 t}$得到
$$
\int_2^{\infty}\left|\frac{E(t)}{t^2}\right|dt \leqslant C\int_2^{\infty} \frac{1}{t\ln^2 t}dt=\frac{C}{\ln 2}
$$
这就说明$\int_2^x \frac{E(t)}{t^2}dt$是绝对收敛的,必然收敛到一个确定的值,我们记为$C_1$.而
$$
\left|\int_x^{\infty}\frac{E(t)}{t^2}dt\right|\leqslant C\int_x^{\infty} \frac{1}{t\ln^2 t}dt = \frac{C}{\ln x}
$$
因此
$$
\begin{aligned}
    \sum_{p\leqslant x}\frac{1}{p} &= \frac{\pi(x)}{x} + \int_2^x \frac{\pi(t)}{t^2}dt\\
    &= \frac{1}{\ln x} + O\left(\frac{1}{\ln^2 x}\right) + \ln\ln x-\ln\ln 2 + C_1 + O\left(\frac{1}{\ln x}\right) \\
\end{aligned}
$$
记常数$M=C_1-\ln\ln 2$,以及合并三个衰减项$\frac{1}{\ln x}+O\left(\frac{1}{\ln^2 x}\right) + O\left(\frac{1}{\ln x}\right) = O\left(\frac{1}{\ln x}\right)$,就得到了Mertens\texttt{'} Second Theorem:
$$
\sum_{p\leqslant x}\frac{1}{p} = \ln\ln x + M + O\left(\frac{1}{\ln x}\right)
$$
你可以看出,所有素数的倒数和是发散的.




\printbibliography
\end{document}
